\thispagestyle{plain}
\section*{底本与版权说明}
\addcontentsline{toc}{section}{底本与版权说明}

\begin{itemize}[leftmargin=0em,itemsep=0.45em,label={}]
\item \textbf{底本}\quad Lin Chen, \emph{Lectures on Infinity Categories}, 2024. 全书二十七讲，正文 203 页（PDF 页）。作者为清华大学求真书院教授，讲义为其“无穷范畴”课程（2024 年秋季）的课堂讲义，未正式出版。
\item \textbf{底本获取}\quad 讲义 PDF 由作者在课程中发布，本译本据该 PDF 全文译出。底本文件不收入本仓库，仅存于仓库外的工作目录。
\item \textbf{页码口径}\quad 底本的印刷页码在每讲内从 1 重新计数，PDF 页则是全书的连续编号（1–203）。本译本每讲开头标注该讲的“底本 PDF 第 $x$–$y$ 页”；资料卡与译注中若标“原书页”，指该讲内的印刷页码，换算关系为：PDF 页 $=$ 该讲起始 PDF 页 $+$ 原书页 $-1$。
\item \textbf{译文性质}\quad 译文是对底本的演绎作品，仅供学术研读、教学与检索使用。如权利人提出异议，本译本将撤下。
\item \textbf{译注}\quad 正文中的竖线小字为译者注，讨论术语取舍、背景补充或原文疑点。凡解释性译注均标注依据等级：\textbf{A} 为可核对的公开文献，\textbf{B} 为同类教材通例，\textbf{I} 为译者据上下文推断；无外部依据者标“译者按（无外部依据）”。
\item \textbf{原文处理}\quad 原文明显的拼写笔误（如 pathes、wihch、difficult 等）与少量符号笔误径行改正，不另出注；影响理解的疑点收入《资料卡/04-术语文献勘误.md》。
\item \textbf{标题}\quad 原书每讲只标“LECTURE $n$”，不设讲题。本译本据该讲内容补出讲题，仅为便览，不属于原文。
\end{itemize}
